\documentclass[11pt]{article}
% Viva preparation: notation, full derivation, and likely questions from a PDE professor.
% Build:  cd prep && latexmk -pdf viva_prep.tex
\usepackage[a4paper,margin=2.3cm]{geometry}
\usepackage{amsmath,amssymb,amsthm}
\usepackage{booktabs}
\usepackage{siunitx}
\usepackage{enumitem}
\usepackage{xcolor}
\usepackage[colorlinks=true,linkcolor=blue!50!black]{hyperref}
\sisetup{exponent-product=\times}

% Auto-generated by scripts/export_macros.py from results/results.json. Do not edit.
\newcommand{\Kval}{\num{100}}
\newcommand{\rval}{\num{0.05}}
\newcommand{\sigmaval}{\num{0.25}}
\newcommand{\Tval}{\num{1}}
\newcommand{\Bone}{\num{80}}
\newcommand{\Btwo}{\num{130}}
\newcommand{\kval}{\num{1.6}}
\newcommand{\alphaval}{\num{-0.3}}
\newcommand{\betaval}{\num{-1.69}}
\newcommand{\taufinal}{\num{0.03125}}
\newcommand{\priceEighty}{\num{3.1415}}
\newcommand{\priceNinety}{\num{6.8698}}
\newcommand{\priceHundred}{\num{12.3360}}
\newcommand{\priceHundredTen}{\num{19.3051}}
\newcommand{\priceHundredTwenty}{\num{27.4063}}
\newcommand{\TableMethodsRows}{80 & \num{3.1415} & \num{3.1415} & \num{3.1398} & \num{3.1403} & \num{3.150} & [\num{3.133}, \num{3.168}] \\
90 & \num{6.8698} & \num{6.8698} & \num{6.8686} & \num{6.8675} & \num{6.888} & [\num{6.861}, \num{6.914}] \\
100 & \num{12.3360} & \num{12.3360} & \num{12.3351} & \num{12.3332} & \num{12.360} & [\num{12.324}, \num{12.396}] \\
110 & \num{19.3051} & \num{19.3051} & \num{19.3043} & \num{19.3026} & \num{19.335} & [\num{19.290}, \num{19.380}] \\
120 & \num{27.4063} & \num{27.4063} & \num{27.4054} & \num{27.4046} & \num{27.441} & [\num{27.387}, \num{27.494}] \\}
\newcommand{\errHeatKernel}{\num{2.5e-14}}
\newcommand{\errExplicit}{\num{0.0017}}
\newcommand{\errCN}{\num{0.0028}}
\newcommand{\errMC}{\num{0.034}}
\newcommand{\mcCoverCount}{5}
\newcommand{\mcSeMax}{\num{0.027}}
\newcommand{\fdDx}{\num{0.0125}}
\newcommand{\fdLam}{\num{0.4}}
\newcommand{\cnSteps}{\num{400}}
\newcommand{\mcPaths}{\num{1e6}}
\newcommand{\parityGap}{\num{5.7e-14}}
\newcommand{\stabDx}{\num{0.0125}}
\newcommand{\stabErrThirty}{\num{0.0016}}
\newcommand{\stabErrForty}{\num{0.0018}}
\newcommand{\stabErrHalf}{\num{0.0066}}
\newcommand{\stabErrFiftyTwo}{\num{2.2e10}}
\newcommand{\stabErrFiftyFive}{\num{2.2e26}}
\newcommand{\stabErrSixty}{\num{5.7e46}}
\newcommand{\TableStabilityRows}{0.3 & \num{0.2999} & 667 & \num{1} & \num{0.0016} \\
0.4 & \num{0.4000} & 500 & \num{1} & \num{0.0018} \\
0.5 & \num{0.5000} & 400 & \num{1} & \num{0.0066} \\
0.52 & \num{0.5195} & 385 & \num{1.08} & \num{2.2e10} \\
0.55 & \num{0.5495} & 364 & \num{1.2} & \num{2.2e26} \\
0.6 & \num{0.6006} & 333 & \num{1.4} & \num{5.7e46} \\}
\newcommand{\convOrderMax}{2.00}
\newcommand{\convOrderK}{2.00}
\newcommand{\convLocalMin}{2.000}
\newcommand{\convLocalMax}{2.002}
\newcommand{\TableConvergenceRows}{\num{0.025} & 125 & \num{0.00741} & \num{0.0035} \\
\num{0.0125} & 500 & \num{0.00185} & \num{8.79e-4} \\
\num{0.00625} & 2000 & \num{4.62e-4} & \num{2.20e-4} \\
\num{0.003125} & 8000 & \num{1.16e-4} & \num{5.50e-5} \\
\num{0.0015625} & 32000 & \num{2.89e-5} & \num{1.38e-5} \\}
\newcommand{\convErrFinest}{\num{2.9e-5}}
\newcommand{\convOrderOff}{2.00}
\newcommand{\convLocalOffMin}{2.00}
\newcommand{\convLocalOffMax}{2.00}
\newcommand{\convErrOffFinest}{\num{2.7e-5}}
\newcommand{\barA}{\num{-0.2231}}
\newcommand{\barB}{\num{0.2624}}
\newcommand{\barEll}{\num{0.4855}}
\newcommand{\barDiffMax}{\num{1.1e-5}}
\newcommand{\TableBarrierRows}{85 & \num{0.7610} & \num{4.7841} \\
90 & \num{1.3657} & \num{6.8698} \\
100 & \num{1.9621} & \num{12.3360} \\
110 & \num{1.7492} & \num{19.3051} \\
120 & \num{0.9755} & \num{27.4063} \\
125 & \num{0.4911} & \num{31.7656} \\}
\newcommand{\barPriceHundred}{\num{1.9621}}
\newcommand{\termsShortest}{51}
\newcommand{\termsLongest}{2}
\newcommand{\termsShortestDays}{1}
\newcommand{\termsPredShortest}{51}
\newcommand{\termsPredLongest}{2.7}
\newcommand{\nnPlainRmseTwenty}{\num{0.78}}
\newcommand{\nnPinnRmseTwenty}{\num{0.021}}
\newcommand{\nnPlainResTwenty}{\num{0.033}}
\newcommand{\nnPinnResTwenty}{\num{6.6e-4}}
\newcommand{\nnRmseRatioTwenty}{37}
\newcommand{\nnResRatioTwenty}{50}
\newcommand{\nnPinnSecondsTwenty}{111}
\newcommand{\nnPlainRmseFifty}{\num{0.31}}
\newcommand{\nnPinnRmseFifty}{\num{0.028}}
\newcommand{\nnPlainResFifty}{\num{0.017}}
\newcommand{\nnPinnResFifty}{\num{7.4e-4}}
\newcommand{\nnRmseRatioFifty}{11}
\newcommand{\nnResRatioFifty}{22}
\newcommand{\nnPinnSecondsFifty}{102}
\newcommand{\nnPlainRmseTwoHundred}{\num{0.2}}
\newcommand{\nnPinnRmseTwoHundred}{\num{0.026}}
\newcommand{\nnPlainResTwoHundred}{\num{0.014}}
\newcommand{\nnPinnResTwoHundred}{\num{7.7e-4}}
\newcommand{\nnRmseRatioTwoHundred}{8}
\newcommand{\nnResRatioTwoHundred}{18}
\newcommand{\nnPinnSecondsTwoHundred}{119}
\newcommand{\TableNNRows}{20 & \num{0.78} $\pm$ \num{0.2} & \num{0.021} $\pm$ \num{0.01} & \num{0.033} & \num{6.6e-4} \\
50 & \num{0.31} $\pm$ \num{0.1} & \num{0.028} $\pm$ \num{0.01} & \num{0.017} & \num{7.4e-4} \\
200 & \num{0.2} $\pm$ \num{0.07} & \num{0.026} $\pm$ \num{0.01} & \num{0.014} & \num{7.7e-4} \\}
\newcommand{\TableNNRowsShort}{20 & \num{0.78} & \num{0.021} & \num{0.033} & \num{6.6e-4} \\
50 & \num{0.31} & \num{0.028} & \num{0.017} & \num{7.4e-4} \\
200 & \num{0.2} & \num{0.026} & \num{0.014} & \num{7.7e-4} \\}
\newcommand{\nnAdamIters}{\num{10000}}
\newcommand{\nnLbfgsIters}{\num{3000}}
\newcommand{\nnColloc}{\num{4096}}
\newcommand{\nnTerminal}{\num{400}}
\newcommand{\nnBoundary}{\num{200}}
\newcommand{\nnSeeds}{3}
\newcommand{\nnWidth}{32}
\newcommand{\nnDepth}{3}
\newcommand{\nnLr}{\num{0.002}}
\newcommand{\nnSmax}{\num{2}}
\newcommand{\nnThetaMin}{\num{0.02}}
\newcommand{\nnTestGrid}{60}
\newcommand{\illEps}{\num{1.00e-6}}
\newcommand{\illModes}{127}
\newcommand{\illMaxNoisyHalfMilli}{\num{1}}
\newcommand{\illMaxNoisyOneMilli}{\num{2.8}}
\newcommand{\illMaxNoisyTwoMilli}{\num{1.3e7}}
\newcommand{\illAmpOneMilli}{\num{1.0e7}}
\newcommand{\illMaxCleanThreeMilli}{\num{1.2e+03}}
\newcommand{\hadTwentyData}{\num{0.05}}
\newcommand{\hadTwentySol}{\num{1.2e16}}

\newcommand{\N}{\mathcal N}
\newcommand{\dd}{\,\mathrm d}
\newcommand{\R}{\mathbb R}

% Question / answer layout
\newcounter{qn}[section]
\renewcommand{\theqn}{\thesection.\arabic{qn}}
\newcommand{\Q}[1]{\par\medskip\refstepcounter{qn}\noindent\textbf{Q\theqn.\ #1}\par\nopagebreak\smallskip}
\newcommand{\A}{\noindent\textit{Answer.}\ }

\theoremstyle{definition}
\newtheorem*{remark}{Remark}

\title{Black--Scholes is the Heat Equation in Disguise\\[4pt]
\large Viva preparation: notation, full derivation, and likely questions}
\author{Harshil Goyal \and Amogh Bindal}
\date{MA5030 Partial Differential Equations}

\begin{document}
\maketitle
\tableofcontents

% =====================================================================
\section{Notation: every symbol in the talk}

\subsection*{Financial variables}
\begin{center}
\renewcommand{\arraystretch}{1.25}
\begin{tabular}{@{}lp{0.72\textwidth}@{}}
\toprule
Symbol & Meaning \\
\midrule
$S$ & Price of the underlying share (stock) at time $t$; $S>0$. Independent variable of the PDE. \\
$t$ & Calendar time, $0\le t\le T$. Independent variable of the PDE. \\
$T$ & Expiry (maturity) date of the option. Default $T=\Tval$ year. \\
$K$ & Strike price: the fixed price at which the option lets you buy. Default $K=\Kval$. \\
$V(S,t)$ & Value (fair price) of the option when the share price is $S$ at time $t$. The unknown function. \\
$\max(S-K,0)$ & Payoff of a European call at expiry, also written $(S-K)^+$. \\
$r$ & Risk-free interest rate (continuously compounded bank rate). Default $r=\rval$. \\
$\mu$ & Drift: expected growth rate of the share. Disappears from the final equation. \\
$\sigma$ & Volatility: size of the random fluctuations of $\ln S$ per $\sqrt{\text{year}}$. Default $\sigma=\sigmaval$. \\
$W_t$ & Standard Brownian motion (Wiener process): $W_0=0$, independent Gaussian increments, $W_{t+h}-W_t\sim N(0,h)$. \\
$\Delta$ & Number of shares held in the hedge; chosen as $\Delta=\partial V/\partial S$ (``the delta''). \\
$\Pi$ & Value of the hedged portfolio $\Pi=V-\Delta S$. \\
$B_1,B_2$ & Lower and upper barriers of the double-barrier option ($\Bone$, $\Btwo$). \\
\bottomrule
\end{tabular}
\end{center}

\subsection*{Transformed (heat-equation) variables}
\begin{center}
\renewcommand{\arraystretch}{1.25}
\begin{tabular}{@{}lp{0.72\textwidth}@{}}
\toprule
Symbol & Meaning \\
\midrule
$x=\ln(S/K)$ & Log-moneyness: new space variable, $x\in\R$. $S=K$ is $x=0$; $S\to0$ is $x\to-\infty$. \\
$\tau=\tfrac12\sigma^2(T-t)$ & Scaled time to expiry, dimensionless. $t=T$ is $\tau=0$. \\
$k=2r/\sigma^2$ & The only parameter left after scaling; here $k=\kval$. \\
$v(x,\tau)$ & $V=K\,v$: price measured in units of the strike. \\
$\alpha,\beta$ & Exponents of the factor $e^{\alpha x+\beta\tau}$: $\alpha=-\frac{k-1}{2}=\alphaval$, $\beta=-\frac{(k+1)^2}{4}=\betaval$. \\
$u(x,\tau)$ & $V=K\,e^{\alpha x+\beta\tau}u$: solves the heat equation $u_\tau=u_{xx}$. \\
$u_0(x)$ & Initial temperature $u(x,0)=\max\big(e^{\frac{k+1}2x}-e^{\frac{k-1}2x},0\big)$. \\
$\Phi(x,\tau)$ & Heat kernel (fundamental solution) $\frac{1}{\sqrt{4\pi\tau}}e^{-x^2/4\tau}$. \\
$\N(d)$ & Standard normal distribution function $\frac{1}{\sqrt{2\pi}}\int_{-\infty}^d e^{-z^2/2}\dd z$. \\
$d_1,d_2$ & Arguments of $\N$ in the Black--Scholes formula. \\
$a,b,\ell$ & $a=\ln(B_1/K)$, $b=\ln(B_2/K)$, $\ell=b-a$: barrier interval in $x$. \\
$\lambda=\Delta\tau/\Delta x^2$ & Mesh ratio of the explicit finite-difference scheme (not the eigenvalue). \\
$\lambda_n=(n\pi/\ell)^2$ & Eigenvalues of $-X''$ on $[a,b]$ with zero boundary values. \\
\bottomrule
\end{tabular}
\end{center}

\noindent Subscripts denote partial derivatives: $V_t=\partial V/\partial t$, $V_{SS}=\partial^2V/\partial S^2$, etc.

% =====================================================================
\section{The derivation from start to end}

\subsection{What is being priced}
A European call gives the holder the \emph{right, not the obligation}, to buy one share for $K$ at time $T$.
At $T$ you exercise only if $S>K$, gaining $S-K$, so
\[
  V(S,T)=\max(S-K,0).
\]
The question is the fair value $V(S,t)$ for $t<T$, when the future price is still random.

\subsection{Model for the share price}
The share follows geometric Brownian motion (GBM)
\[
  \dd S=\mu S\dd t+\sigma S\dd W .
\]
The return $\dd S/S$ has a deterministic part $\mu\dd t$ and a random part $\sigma\dd W$ with
$\dd W\sim N(0,\dd t)$. The solution is $S_t=S_0\exp\big((\mu-\tfrac12\sigma^2)t+\sigma W_t\big)$, so
$S>0$ always and $\ln S$ is normally distributed (lognormal prices).

\subsection{Itô's lemma (the chain rule with noise)}
Taylor-expand $V(S,t)$:
\[
  \dd V=V_t\dd t+V_S\dd S+\tfrac12V_{SS}(\dd S)^2+V_{St}\dd S\dd t+\tfrac12V_{tt}(\dd t)^2+\cdots
\]
In ordinary calculus $(\dd S)^2$ is negligible. Here $\dd W$ is of size $\sqrt{\dd t}$, and the rules
\[
  (\dd W)^2=\dd t,\qquad \dd W\dd t=0,\qquad (\dd t)^2=0
\]
give $(\dd S)^2=\sigma^2S^2(\dd W)^2=\sigma^2S^2\dd t$, which is first order and must be kept. Hence
\[
  \dd V=\Big(V_t+\mu SV_S+\tfrac12\sigma^2S^2V_{SS}\Big)\dd t+\sigma SV_S\dd W. \tag{1}
\]

\subsection{Delta hedging and the PDE}
Hold one option and sell $\Delta$ shares: $\Pi=V-\Delta S$. Over $\dd t$ (with $\Delta$ held fixed),
\[
  \dd\Pi=\dd V-\Delta\dd S
  =\Big(V_t+\mu SV_S+\tfrac12\sigma^2S^2V_{SS}-\mu\Delta S\Big)\dd t+\sigma S\,(V_S-\Delta)\dd W .
\]
Choose $\Delta=V_S$. The $\dd W$ term vanishes, so the portfolio is riskless, and the $\mu$ terms cancel too:
\[
  \dd\Pi=\Big(V_t+\tfrac12\sigma^2S^2V_{SS}\Big)\dd t .
\]
\textbf{No arbitrage:} a riskless portfolio must earn exactly the bank rate, otherwise one could borrow at
$r$ and invest in $\Pi$ (or the reverse) for a risk-free profit. So
$\dd\Pi=r\Pi\dd t=r(V-SV_S)\dd t$. Equating,
\[
  \boxed{\;V_t+\tfrac12\sigma^2S^2V_{SS}+rSV_S-rV=0,\qquad S>0,\ t<T.\;} \tag{BS}
\]
Meaning of the terms: $V_t$ time decay, $\frac12\sigma^2S^2V_{SS}$ diffusion (randomness),
$rSV_S$ drift (growth of money), $-rV$ discounting.

\subsection{Terminal and boundary conditions}
\begin{itemize}[nosep]
  \item $V(S,T)=\max(S-K,0)$ (terminal condition).
  \item $S=0$: GBM started at $0$ stays at $0$, so the option is worthless: $V(0,t)=0$. This is not an extra
    assumption: at $S=0$ (BS) reduces to the ODE $V_t=rV$ with $V(0,T)=0$, so $V(0,t)=0$.
  \item $S\to\infty$: the option is almost surely exercised, so it behaves like the share minus the discounted
    strike: $V\sim S-Ke^{-r(T-t)}$.
\end{itemize}

\subsection{Classification}
For $A w_{XX}+2Bw_{XY}+Cw_{YY}+\text{(lower order)}=0$ the type is decided by $B^2-AC$
($>0$ hyperbolic, $=0$ parabolic, $<0$ elliptic), and it is invariant under smooth invertible changes of
variables. For (BS) with $(X,Y)=(S,t)$: $A=\tfrac12\sigma^2S^2$, $B=C=0$, so $B^2-AC=0$:
\emph{parabolic} for every $S>0$. Characteristics $A\dd t^2=0$ are the lines $t=\text{const}$.
At $S=0$ all of $A,B,C$ vanish (\emph{degenerate}). Since $V_t$ and $V_{SS}$ have the same sign
(compare $u_\tau-u_{xx}=0$ which has opposite signs), it is a \emph{backward} parabolic equation: well posed
when solved from $t=T$ backwards.

\subsection{Step 1: logarithmic price}
Put $S=Ke^x$, i.e.\ $x=\ln(S/K)$, and $V(S,t)=K\,v(x,\tau)$. Since $\frac{\dd x}{\dd S}=\frac1S$,
\[
  V_S=\frac{K}{S}v_x\ \Rightarrow\ SV_S=Kv_x .
\]
Differentiate $SV_S=Kv_x$ with respect to $S$ (product rule on the left, chain rule on the right):
\[
  V_S+SV_{SS}=\frac{K}{S}v_{xx}\ \Rightarrow\ S^2V_{SS}=Kv_{xx}-SV_S=K(v_{xx}-v_x).
\]
The variable coefficients disappear because $S\partial_S=\partial_x$ (Euler--Cauchy type operator).
The half-line $S\in(0,\infty)$ becomes the whole line $x\in\R$.

\subsection{Step 2: reverse and rescale time}
Put $\tau=\tfrac12\sigma^2(T-t)$. Then $\partial\tau/\partial t=-\tfrac12\sigma^2$, so $V_t=-\tfrac12\sigma^2Kv_\tau$.
Substituting all three derivatives into (BS) and dividing by $K$:
\[
  -\tfrac{\sigma^2}{2}v_\tau+\tfrac{\sigma^2}{2}(v_{xx}-v_x)+rv_x-rv=0 .
\]
Divide by $\sigma^2/2$ and set $k=2r/\sigma^2$:
\[
  v_\tau=v_{xx}+(k-1)v_x-kv,\qquad v(x,0)=\frac{1}{K}\max(Ke^x-K,0)=\max(e^x-1,0).
\]
Now $\tau=0$ corresponds to $t=T$: the \emph{terminal} condition has become an \emph{initial} condition, and
the equation runs forward in $\tau$. Four parameters $(K,r,\sigma,T)$ reduced to one, $k$.

\subsection{Step 3: remove the first- and zeroth-order terms}
Try $v=e^{\alpha x+\beta\tau}u$ and write $E=e^{\alpha x+\beta\tau}$. Then
\[
  v_\tau=E(\beta u+u_\tau),\quad v_x=E(\alpha u+u_x),\quad v_{xx}=E(\alpha^2u+2\alpha u_x+u_{xx}).
\]
Substitute and divide by $E>0$:
\[
  \beta u+u_\tau=\alpha^2u+2\alpha u_x+u_{xx}+(k-1)(\alpha u+u_x)-ku,
\]
\[
  u_\tau=u_{xx}+(2\alpha+k-1)u_x+\big(\alpha^2+(k-1)\alpha-k-\beta\big)u .
\]
\begin{itemize}[nosep]
  \item $u_x$ term vanishes: $\alpha=-\frac{k-1}{2}$.
  \item Then $\alpha^2+(k-1)\alpha-k=\frac{(k-1)^2}{4}-\frac{(k-1)^2}{2}-k=-\frac{(k-1)^2+4k}{4}=-\frac{(k+1)^2}{4}$,
  so the $u$ term vanishes for $\beta=-\frac{(k+1)^2}{4}$.
\end{itemize}
Initial data: $u_0(x)=e^{-\alpha x}\max(e^x-1,0)=e^{\frac{k-1}{2}x}\max(e^x-1,0)
=\max\big(e^{\frac{k+1}{2}x}-e^{\frac{k-1}{2}x},0\big)$.

\begin{remark}
Equivalent two-move version: a moving frame $y=x+(k-1)\tau$ removes the advection term (Galilean shift), and
$e^{-k\tau}$ removes the discounting. The single exponential factor does both at once.
\end{remark}

\subsection{Result}
\[
  u_\tau=u_{xx},\quad x\in\R,\ \tau>0,\qquad u(x,0)=u_0(x),\qquad
  V(S,t)=K\,e^{\alpha x+\beta\tau}u(x,\tau).
\]
Each substitution is smooth and invertible, so solving the heat problem is the same as solving (BS).

% ---------------------------------------------------------------------
\subsection{Solving the heat equation: separation of variables}
Try $u=X(x)T(\tau)$. Then $XT'=X''T$, so $\frac{T'}{T}=\frac{X''}{X}=$ const.
\begin{itemize}[nosep]
  \item Constant $=+\kappa^2>0$: $X=e^{\pm\kappa x}$, unbounded on $\R$: rejected.
  \item Constant $=-\mu^2\le0$: $T=e^{-\mu^2\tau}$, $X=A\cos\mu x+B\sin\mu x$ (bounded).
\end{itemize}
With no boundary conditions every $\mu\ge0$ is allowed, so we superpose with an integral instead of a sum:
\[
  u(x,\tau)=\int_0^\infty e^{-\mu^2\tau}\big(A(\mu)\cos\mu x+B(\mu)\sin\mu x\big)\dd\mu .
\]
At $\tau=0$ this is the Fourier integral of $g=u_0$, so
\[
  A(\mu)=\frac1\pi\int_\R g(y)\cos\mu y\dd y,\qquad B(\mu)=\frac1\pi\int_\R g(y)\sin\mu y\dd y .
\]
Substituting and using $\cos\mu y\cos\mu x+\sin\mu y\sin\mu x=\cos\mu(y-x)$, then swapping the order of integration,
\[
  u(x,\tau)=\frac1\pi\int_\R g(y)\,I\dd y,\qquad I=\int_0^\infty e^{-\mu^2\tau}\cos\mu(y-x)\dd\mu .
\]

\subsection{Evaluating $I$}
Let $z=\mu\sqrt\tau$, $\eta=(y-x)/\sqrt\tau$, so $I=\frac{1}{\sqrt\tau}K(\eta)$ with
$K(\eta)=\int_0^\infty e^{-z^2}\cos\eta z\dd z$. Differentiate under the integral and integrate by parts with
$ze^{-z^2}=-\frac12\frac{\dd}{\dd z}e^{-z^2}$:
\[
  K'(\eta)=-\int_0^\infty ze^{-z^2}\sin\eta z\dd z
  =\Big[\tfrac12e^{-z^2}\sin\eta z\Big]_0^\infty-\frac\eta2\int_0^\infty e^{-z^2}\cos\eta z\dd z
  =-\frac\eta2K(\eta).
\]
So $K(\eta)=K(0)e^{-\eta^2/4}$ with $K(0)=\int_0^\infty e^{-z^2}\dd z=\frac{\sqrt\pi}{2}$. Therefore
$I=\frac12\sqrt{\frac{\pi}{\tau}}\,e^{-(y-x)^2/4\tau}$ and
\[
  \boxed{\;u(x,\tau)=\int_\R\Phi(x-y,\tau)\,u_0(y)\dd y,\qquad \Phi(x,\tau)=\frac{1}{\sqrt{4\pi\tau}}e^{-x^2/4\tau}.\;}
\]
(Second route, by scaling: $u_\tau=u_{xx}$ is invariant under $x\to\lambda x$, $\tau\to\lambda^2\tau$, so look
for $u=\tau^{-1/2}v(x^2/\tau)$; this gives $v=e^{-z/4}$ and the same $\Phi$ after normalising
$\int\Phi\dd x=1$.)

\paragraph{Properties of $\Phi$.} $\Phi>0$; smooth for $\tau>0$; solves the heat equation; $\int_\R\Phi\dd x=1$;
as $\tau\to0^+$ it concentrates at $x=0$ (it tends to the Dirac delta).

\paragraph{Theorem (existence).} If $u_0$ is continuous with $|u_0(y)|\le Me^{c|y|}$, then $u=\Phi*u_0$ is
$C^\infty$ for $\tau>0$, solves $u_\tau=u_{xx}$, and $u(x,\tau)\to u_0(x_0)$ as $(x,\tau)\to(x_0,0^+)$.
\emph{Proof idea:} differentiate under the integral (Gaussian decay dominates $e^{c|y|}$); for the initial
value, split $\int\Phi\,|u_0(y)-u_0(x_0)|\dd y$ into $|y-x_0|<\delta$ (small by continuity, since $\int\Phi=1$)
and $|y-x_0|\ge\delta$ (small because $\Phi\to0$ uniformly away from $x$).

\paragraph{Uniqueness.} Among solutions with $|u|\le Ae^{ax^2}$ the solution is unique (maximum principle on
the whole line). Our $u$ grows only like $e^{c|x|}$, so the price is unique.

\subsection{Computing the convolution}\label{sec:conv}
Since $u_0(s)=0$ for $s<0$,
\[
  u=I_+-I_-,\qquad I_\pm=\frac{1}{\sqrt{4\pi\tau}}\int_0^\infty e^{a_\pm s}e^{-(x-s)^2/4\tau}\dd s,\quad a_\pm=\tfrac{k\pm1}{2}.
\]
Substitute $s=x+\sqrt{2\tau}\,y$: $\dd s=\sqrt{2\tau}\dd y$, $\frac{(x-s)^2}{4\tau}=\frac{y^2}{2}$, the lower limit
$s=0$ becomes $y=-x/\sqrt{2\tau}$, and $\frac{\sqrt{2\tau}}{\sqrt{4\pi\tau}}=\frac{1}{\sqrt{2\pi}}$. With
$m_\pm=a_\pm\sqrt{2\tau}$,
\[
  I_\pm=\frac{e^{a_\pm x}}{\sqrt{2\pi}}\int_{-x/\sqrt{2\tau}}^\infty e^{-\frac12y^2+m_\pm y}\dd y .
\]
Complete the square: $-\frac12y^2+m y=-\frac12(y-m)^2+\frac12m^2$. With $z=y-m$ and the symmetry of the
Gaussian,
\[
  I_\pm=e^{a_\pm x+\frac12m_\pm^2}\,\frac{1}{\sqrt{2\pi}}\int_{-d_\pm}^\infty e^{-z^2/2}\dd z
  =e^{a_\pm x+a_\pm^2\tau}\,\N(d_\pm),\qquad d_\pm=\frac{x}{\sqrt{2\tau}}+a_\pm\sqrt{2\tau}.
\]

\subsection{Back to the original variables}
$V=Ke^{\alpha x+\beta\tau}(I_+-I_-)$. Collect the exponents:
\begin{itemize}[nosep]
  \item $I_+$: $\alpha+a_+=-\frac{k-1}2+\frac{k+1}2=1$ and $\beta+a_+^2=0$, so the prefactor is $Ke^x=S$.
  \item $I_-$: $\alpha+a_-=0$ and $\beta+a_-^2=\frac{(k-1)^2-(k+1)^2}{4}=-k$, so the prefactor is
    $Ke^{-k\tau}=Ke^{-r(T-t)}$ because $k\tau=\frac{2r}{\sigma^2}\cdot\frac{\sigma^2}{2}(T-t)=r(T-t)$.
  \item $\sqrt{2\tau}=\sigma\sqrt{T-t}$ and $x=\ln(S/K)$, so
    $d_\pm=\frac{\ln(S/K)}{\sigma\sqrt{T-t}}+\big(\frac{r}{\sigma^2}\pm\frac12\big)\sigma\sqrt{T-t}$.
\end{itemize}
\[
  \boxed{\;V(S,t)=S\,\N(d_1)-Ke^{-r(T-t)}\N(d_2),\qquad
  d_{1,2}=\frac{\ln\frac SK+\big(r\pm\frac{\sigma^2}{2}\big)(T-t)}{\sigma\sqrt{T-t}}.\;}
\]
For the default parameters $V(100,0)=\priceHundred$.

\subsection{Double barrier: Fourier sine series}
If the option dies when $S$ hits $B_1$ or $B_2$, then $u(a,\tau)=u(b,\tau)=0$ on the finite interval $[a,b]$.
Separation gives $X''+\lambda X=0$, $X(a)=X(b)=0$: $\lambda\le0$ gives only $X\equiv0$, and
$\lambda_n=(n\pi/\ell)^2$, $X_n=\sin\frac{n\pi(x-a)}{\ell}$. Orthogonality
$\int_a^bX_mX_n\dd x=\frac\ell2\delta_{mn}$ gives
\[
  u=\sum_{n\ge1}c_n X_n(x)e^{-\lambda_n\tau},\qquad c_n=\frac2\ell\int_{\max(a,0)}^bu_0(x)X_n(x)\dd x .
\]

\subsection{Explicit finite differences}
$u_j^{m+1}=u_j^m+\lambda(u_{j+1}^m-2u_j^m+u_{j-1}^m)$ with $\lambda=\Delta\tau/\Delta x^2$. Inserting the Fourier mode
$u_j^m=G^me^{i\xi j\Delta x}$ gives the amplification factor $G=1-4\lambda\sin^2\frac{\xi\Delta x}{2}\in[1-4\lambda,1]$,
and $|G|\le1$ for all $\xi$ iff $1-4\lambda\ge-1$, i.e.\ $\lambda\le\frac12$. The truncation error is
$O(\Delta\tau+\Delta x^2)=O(\Delta x^2)$ at fixed $\lambda$: second order in space.

% =====================================================================
\section{Likely questions and answers}

\subsection{Modelling and the derivation}

\Q{What exactly is $V$? Is it a function or a random variable?}
\A $V(S,t)$ is a deterministic function of two real variables. The randomness is in $S_t$; the option's
value at time $t$ is the random number $V(S_t,t)$, but $V$ itself is an ordinary function solving a PDE.

\Q{Why does $\mu$ not appear in the equation? Surely a share expected to rise makes the call worth more.}
\A With $\Delta=V_S$ the $\mu$ terms cancel exactly. The option can be replicated by trading the share and
the bank account, so its price is the cost of replication, which does not depend on anyone's view of the
drift. The expected growth is already reflected in today's share price $S$.

\Q{Why must the riskless portfolio earn exactly $r$?}
\A No arbitrage. If it earned more than $r$, borrow at $r$, buy $\Pi$ and pocket a riskless profit; if less,
short $\Pi$ and lend at $r$. Either way prices would adjust until the return is $r$.

\Q{You held $\Delta$ fixed over $\dd t$ even though $\Delta=V_S$ changes. Is that legitimate?}
\A It is the self-financing condition: $\Delta$ is set at the start of each interval and the portfolio
changes only through price moves. A rigorous version builds a self-financing replicating portfolio and
uses the martingale representation theorem; it gives the same PDE.

\Q{What is Itô's lemma, and why is there an extra term compared with the ordinary chain rule?}
\A It is the chain rule for functions of a diffusion. Brownian increments are of size $\sqrt{\dd t}$, so
$(\dd W)^2=\dd t$ is first order. The second-order Taylor term $\frac12V_{SS}(\dd S)^2=\frac12\sigma^2S^2V_{SS}\dd t$
survives. That term is the diffusion term of the PDE.

\Q{Why geometric Brownian motion and not $\dd S=\mu\dd t+\sigma\dd W$?}
\A GBM keeps $S>0$ and makes returns, not absolute changes, scale-free. Arithmetic Brownian motion allows
negative prices. GBM is also what makes the coefficients $S^2$ and $S$ appear, which the logarithm removes.

\Q{Is the Black--Scholes equation linear? What order?}
\A Linear, second order in $S$, first order in $t$, with variable coefficients. Linearity is why
superposition works (e.g.\ the payoff splits into two exponentials) and why put--call parity holds.

\Q{Derive put--call parity from the PDE.}
\A Let $C,P$ be the call and put prices. $W=C-P$ solves (BS) with terminal data
$\max(S-K,0)-\max(K-S,0)=S-K$. The function $S-Ke^{-r(T-t)}$ solves (BS) (check: $V_t=-rKe^{-r(T-t)}$,
$rSV_S-rV=rKe^{-r(T-t)}$, $V_{SS}=0$) with the same terminal data. By uniqueness,
$C-P=S-Ke^{-r(T-t)}$. Numerically our gap is $\le\parityGap$.

\Q{Where does the boundary condition at $S=0$ come from? Do you need to impose it?}
\A No. At $S=0$ the diffusion and drift coefficients both vanish, so (BS) becomes $V_t=rV$ with
$V(0,T)=0$, forcing $V(0,t)=0$. Financially: a share at price $0$ stays there, so the call is worthless.
This is typical of degenerate equations: the degeneracy replaces a boundary condition.

\Q{Why $V\sim S-Ke^{-r(T-t)}$ as $S\to\infty$?}
\A For very large $S$ the option is almost surely exercised, so it is equivalent to receiving the share
and paying $K$ at $T$, worth $S-Ke^{-r(T-t)}$ today. After the change of variables this is the growth
condition at $x\to+\infty$.

\subsection{Classification and well-posedness}

\Q{Classify the equation. Is it parabolic everywhere?}
\A $A=\frac12\sigma^2S^2$, $B=C=0$, so $B^2-AC=0$: parabolic for all $S>0$; degenerate at $S=0$ where $A=0$.
Characteristics are $t=$ const.

\Q{What does ``backward parabolic'' mean, and what goes wrong if you solve forward in $t$?}
\A Writing it as $V_t=-\frac12\sigma^2S^2V_{SS}+\dots$, the diffusion has the ``wrong'' sign for forward
time; forward in $t$ it behaves like $w_t=-w_{xx}$. Hadamard's example: $w(x,0)=\frac1n\sin nx\to0$ but
$w=\frac1ne^{n^2t}\sin nx\to\infty$ for any $t>0$: no continuous dependence on data. That is why the data
are given at $T$ and we solve towards today. Our experiment: noise of $\illEps$ grows to
$\illMaxNoisyTwoMilli$ after time $0.002$.

\Q{Is the problem well posed in the sense of Hadamard?}
\A Yes, backwards in $t$ (forward in $\tau$), within the growth class $|u|\le Ae^{ax^2}$: existence (heat
kernel), uniqueness (maximum principle), and continuous dependence (maximum principle gives
$\sup|u_1-u_2|\le\sup|g_1-g_2|$ for bounded data).

\Q{The type is invariant under changes of variables. Show it.}
\A Under $(X,Y)\to(\xi,\eta)$ with Jacobian $J$, the new coefficients satisfy
$B'^2-A'C'=(\det J)^2(B^2-AC)$. Since $\det J\ne0$ for an invertible map, the sign is preserved.

\Q{What happens if $\sigma\to0$?}
\A The second-order term disappears and (BS) becomes the first-order hyperbolic transport equation
$V_t+rSV_S-rV=0$, solved by characteristics $S(t)=S_0e^{rt}$. The price becomes
$\max(S-Ke^{-r(T-t)},0)$, which keeps the kink: no smoothing without diffusion.

\subsection{The change of variables}

\Q{Why the logarithm?}
\A Because $S\partial_S$ and $S^2\partial_{SS}$ are Euler--Cauchy operators, and $x=\ln S$ turns
$S\partial_S$ into $\partial_x$, giving constant coefficients. It also maps $(0,\infty)$ to $\R$, so the
problem lives on the whole line.

\Q{Show the computation $S^2V_{SS}=K(v_{xx}-v_x)$.}
\A $V_S=\frac{K}{S}v_x$. Differentiating again, $V_{SS}=-\frac{K}{S^2}v_x+\frac{K}{S}\cdot\frac1Sv_{xx}$, so
$S^2V_{SS}=K(v_{xx}-v_x)$. The extra $-v_x$ comes from differentiating $1/S$.

\Q{Why is time reversed and scaled by $\sigma^2/2$?}
\A Reversing ($T-t$) turns the terminal condition into an initial condition and the backward equation into a
forward one. Scaling by $\sigma^2/2$ makes the coefficient of $v_{xx}$ exactly $1$ and $\tau$ dimensionless.

\Q{What is the meaning of $k=2r/\sigma^2$?}
\A The ratio of interest-rate drift to diffusion strength; it is the only parameter left. Here $k=\kval$.
If $k=1$ the first-order term $(k-1)v_x$ is already absent.

\Q{How did you choose $\alpha$ and $\beta$?}
\A Substitute $v=e^{\alpha x+\beta\tau}u$ and require the coefficient of $u_x$ ($2\alpha+k-1$) and of $u$
($\alpha^2+(k-1)\alpha-k-\beta$) to vanish. Two equations, two unknowns:
$\alpha=-\frac{k-1}{2}=\alphaval$, $\beta=-\frac{(k+1)^2}{4}=\betaval$.

\Q{Does this trick work for any linear constant-coefficient parabolic equation?}
\A Yes: $u_\tau=u_{xx}+bu_x+cu$ becomes the heat equation with $\alpha=-b/2$, $\beta=c-b^2/4$. It fails if
coefficients depend on $x$ or $\tau$ in a general way (e.g.\ local volatility $\sigma(S,t)$).

\Q{Are the two problems really equivalent?}
\A Each map is a smooth bijection: $S\leftrightarrow x$, $t\leftrightarrow\tau$, and multiplication by
$Ke^{\alpha x+\beta\tau}>0$. A solution of one problem gives a solution of the other, and uniqueness
transfers.

\subsection{Solving the heat equation}

\Q{In separation of variables, why must the separation constant be negative?}
\A A positive constant gives $X=e^{\pm\kappa x}$, unbounded as $|x|\to\infty$ (and $T$ grows in time). Zero
gives $X=A+Bx$; only the constant is bounded, which is the $\mu=0$ case. So we take $-\mu^2$.

\Q{Why an integral over $\mu$ rather than a sum?}
\A On a bounded interval the boundary conditions select discrete frequencies $n\pi/l$. On $\R$ there are
none, so all $\mu\ge0$ are admissible; as $l\to\infty$ the spacing $\pi/l\to0$ and the Fourier series becomes
the Fourier integral.

\Q{Our $u_0$ grows exponentially, so its Fourier integral does not exist. Is the derivation valid?}
\A The Fourier step is formal for our data. But the final formula $u=\Phi*u_0$ is then proved directly
(existence theorem): the Gaussian decay $e^{-y^2/4\tau}$ beats $e^{c|y|}$, so the convolution converges,
solves the equation and attains the initial data.

\Q{Why can you swap the order of integration?}
\A For $\tau>0$ and integrable $g$, $|e^{-\mu^2\tau}g(y)|$ is integrable over $(\mu,y)$, so Fubini applies.
For growing data, see the previous answer.

\Q{Derive the heat kernel another way.}
\A (i) Scaling/similarity: $u=\tau^{-1/2}v(x^2/\tau)$ gives $4v''+v'=0$, so $v=e^{-z/4}$. (ii) Fourier
transform: $\hat u_\tau=-\xi^2\hat u$, so $\hat u=e^{-\xi^2\tau}\hat g$, and the inverse transform of a Gaussian is
a Gaussian. (iii) Probabilistically, $\Phi$ is the density of $\sqrt2\,W_\tau$.

\Q{Why the factor $1/\sqrt{4\pi\tau}$?}
\A So that $\int_\R\Phi\dd x=1$ for every $\tau$ (conservation of heat). With $x=2\sqrt\tau z$,
$\int e^{-x^2/4\tau}\dd x=2\sqrt\tau\sqrt\pi$.

\Q{What does $\Phi$ look like as $\tau\to0^+$?}
\A Height $\to\infty$ at $x=0$, value $\to0$ elsewhere, total mass $1$: it tends to the Dirac delta. That is why
$\Phi*g\to g$.

\Q{Is the solution of the heat equation on $\R$ unique?}
\A Not without a growth condition: Tychonoff constructed a nonzero solution with zero initial data that
grows faster than $e^{ax^2}$. In the class $|u|\le Ae^{ax^2}$ it is unique by the maximum principle. Our
solution grows like $e^{c|x|}$, so it is unique.

\Q{State and prove the weak maximum principle.}
\A On $\Omega_T=\Omega\times(0,T]$, a solution attains its maximum on the parabolic boundary (bottom and sides).
Proof: $u^\varepsilon=u-\varepsilon\tau$ satisfies $u^\varepsilon_\tau-u^\varepsilon_{xx}=-\varepsilon<0$. At an interior or top maximum
$u^\varepsilon_{xx}\le0$ and $u^\varepsilon_\tau\ge0$, so $u^\varepsilon_\tau-u^\varepsilon_{xx}\ge0$, a contradiction. Let
$\varepsilon\to0$.

\Q{Prove uniqueness on a bounded interval by the energy method.}
\A For the difference $w$ of two solutions, $E(\tau)=\int w^2\dd x$ satisfies $E'=2\int ww_{xx}\dd x=-2\int w_x^2\dd x\le0$
(boundary term vanishes since $w=0$ there). $E(0)=0$ and $E\ge0$, so $E\equiv0$.

\Q{What does the maximum principle say about option prices?}
\A Comparison: if two payoffs satisfy $g_1\le g_2$, their prices satisfy $V_1\le V_2$ (the factor
$Ke^{\alpha x+\beta\tau}$ is positive). In particular payoff $\ge0$ implies price $\ge0$, and $V\le S$.

\Q{What is ``infinite speed of propagation'' here, and what does it mean financially?}
\A Since $\Phi>0$ everywhere, $u(x,\tau)>0$ for all $x$ as soon as $\tau>0$, even though $u_0=0$ for $x<0$. A call
on a share far below the strike is still worth something: there is always some probability of finishing in
the money.

\Q{Why is the price smooth although the payoff has a kink?}
\A Heat-equation smoothing: $\Phi*u_0$ is $C^\infty$ for $\tau>0$. The kink at $S=K$ only exists at $t=T$. As
$t\to T$ the second derivative (``Gamma'') near $S=K$ blows up like $(T-t)^{-1/2}$.

\Q{Walk us through completing the square.}
\A See Section~\ref{sec:conv}. The key steps are the substitution $s=x+\sqrt{2\tau}\,y$ that makes the exponent
$-y^2/2$, then $-\frac12y^2+my=-\frac12(y-m)^2+\frac12m^2$, and the remaining integral is a normal
probability $\N(d_\pm)$.

\Q{Is there a probabilistic interpretation of the formula?}
\A Feynman--Kac: $V(S,t)=e^{-r(T-t)}\mathbb E^{Q}[\max(S_T-K,0)\mid S_t=S]$, where under the risk-neutral
measure $Q$ the drift $\mu$ is replaced by $r$. $\N(d_2)$ is the risk-neutral probability of exercise, and
$\N(d_1)$ is the hedge ratio $\Delta=V_S$. The heat kernel is the transition density of Brownian motion.
Our Monte Carlo check uses exactly this expectation.

\subsection{Barrier option, numerics and neural networks}

\Q{Why does the barrier problem give a sine series?}
\A Knock-out at $B_1,B_2$ means $V=0$ there, so $u(a,\tau)=u(b,\tau)=0$: Dirichlet conditions on a finite
interval. The Sturm--Liouville problem $X''+\lambda X=0$, $X(a)=X(b)=0$ has eigenfunctions
$\sin\frac{n\pi(x-a)}{\ell}$.

\Q{The initial data is nonzero at the upper barrier $x=b$ but the boundary value is $0$. Is that a problem?}
\A The data are incompatible at the corner $(b,0)$. The sine series converges in $L^2$ at $\tau=0$ (with Gibbs
oscillations near $b$), but for every $\tau>0$ the factor $e^{-\lambda_n\tau}$ makes the series converge rapidly and
the solution is smooth inside.

\Q{Why are more terms needed close to expiry?}
\A Mode $n$ is damped by $e^{-(n\pi/\ell)^2\tau}$. For accuracy $\epsilon$ we need $(N\pi/\ell)^2\tau\gtrsim\ln(1/\epsilon)$,
so $N\propto\tau^{-1/2}\propto(T-t)^{-1/2}$.

\Q{Derive the stability condition of the explicit scheme.}
\A Von Neumann: $u_j^m=G^me^{i\xi j\Delta x}$ gives $G=1-4\lambda\sin^2\frac{\xi\Delta x}{2}$. The worst mode is the
sawtooth $\xi\Delta x=\pi$, $G=1-4\lambda$. $|G|\le1$ iff $\lambda\le\frac12$. Above that the sawtooth grows
geometrically: error $\stabErrForty$ at $\lambda=0.4$ but $\stabErrFiftyTwo$ at $\lambda=0.52$.

\Q{Why is Crank--Nicolson stable for every $\lambda$?}
\A Its amplification factor is $G=\frac{1-2\lambda s}{1+2\lambda s}$ with $s=\sin^2\frac{\xi\Delta x}{2}\ge0$, so
$|G|\le1$ always. It is second order in both $\Delta\tau$ and $\Delta x$.

\Q{The payoff has a kink. Why do you still see second-order convergence?}
\A The kink's Fourier coefficients decay like $\xi^{-2}$, and the scheme damps the high modes that carry the
local error, so the global error at positive time is still $O(\Delta x^2)$. Observed order: $\convOrderMax$.

\Q{On a computer you cannot use all of $\R$. What boundary conditions did you use?}
\A Truncate to $[x_{\min},x_{\max}]$ and impose the asymptotic values: $u\approx0$ far left, and the
transformed $S-Ke^{-r(T-t)}$ far right. The error from truncation is exponentially small if the domain is
wide enough.

\Q{What is a PINN and why does it beat the plain network?}
\A A physics-informed neural network adds $\|V_t+\frac12\sigma^2S^2V_{SS}+rSV_S-rV\|^2$ at many collocation
points to the loss, computing derivatives by automatic differentiation, plus terminal and boundary terms.
The PDE supplies information between the few data points, so with $20$ labels it is about
$\nnRmseRatioTwenty$ times more accurate.

\subsection{Limitations and extensions}

\Q{What assumptions break in real markets?}
\A Constant $\sigma$ and $r$, continuous paths (no jumps), no transaction costs, continuous hedging, no
dividends. Real implied volatility depends on $K$ and $T$ (the volatility smile).

\Q{How do dividends change things?}
\A A continuous dividend yield $q$ changes the drift term to $(r-q)SV_S$. The reduction still works with
$k$ replaced appropriately, and $S$ in the formula becomes $Se^{-q(T-t)}$.

\Q{What about American options?}
\A Early exercise gives the constraint $V\ge\max(S-K,0)$ and a free boundary (optimal exercise price). The
problem becomes an obstacle problem / linear complementarity problem; in heat variables it is the heat
equation with an obstacle. There is no closed form for the American put.

\Q{What if $\sigma$ depends on $S$ and $t$?}
\A The log change still works, but the coefficients remain variable, so no exponential factor reduces it to
the heat equation. One solves it numerically (finite differences, Monte Carlo).

\Q{What if the share price jumps?}
\A A jump-diffusion model (e.g.\ Merton) adds an integral term, giving a partial integro-differential equation.
It is no longer a local PDE and the heat-kernel solution does not apply directly.

\subsection{Quick-fire questions}
\begin{itemize}
  \item \emph{Units of $\tau$?} Dimensionless ($\sigma^2$ has units 1/time).
  \item \emph{Value of a call at $S=K$, $t\to T$?} $\to0$, like $S\sigma\sqrt{(T-t)/(2\pi)}$.
  \item \emph{Is $V$ increasing in $\sigma$?} Yes: Vega $=S\sqrt{T-t}\,\N'(d_1)>0$; more diffusion spreads more
    mass above $K$.
  \item \emph{Is $V$ convex in $S$?} Yes, Gamma $V_{SS}=\N'(d_1)/(S\sigma\sqrt{T-t})>0$; convexity of the payoff
    is preserved by the heat flow.
  \item \emph{Bounds?} $\max(S-Ke^{-r(T-t)},0)\le V\le S$.
  \item \emph{Which class tools did you use?} Classification, change of variables to canonical form, separation
    of variables, Fourier integral and series, fundamental solution, maximum principle, energy method,
    Hadamard ill-posedness, von Neumann stability.
\end{itemize}

\end{document}
